\documentclass[reqno]{amsart}
\usepackage[T1]{fontenc}
\usepackage{amsmath,amssymb,amsthm}
\usepackage{xr-hyper}
\usepackage[hidelinks]{hyperref}
\externaldocument[c1-]{20261125T001000--zeros-to-ramp-sums__workbook}
\externaldocument[c3-]{20261125T005000--galois-fourier-and-averaging__workbook}
\externaldocument[c5-]{20261125T000000--slope-method__workbook}
\externaldocument[c6-]{20261126T030100--a-twist-for-each-prime__workbook}
\externaldocument[c32-]{20261126T030300--the-data__workbook}
\externaldocument[c33-]{20261126T020210--bessel-poincare-series-and-their-gram-operators__workbook}
\externaldocument[c36-]{20261126T030600--theorem-a-short-forms__workbook}
\externaldocument[c37-]{20261126T030700--the-determinant-of-a-chain__workbook}
\externaldocument[c38-]{20261126T030800--the-cost-of-one-lift__workbook}
\externaldocument[c40-]{20261126T031000--twists-with-few-congruences__workbook}

% Chapter number; \ref{chap} gives it.
\newcounter{chap}
% Chapter numbers, generated from 20261125T012000--grh-workbook__guide.org
% by scripts/make-chapters.py.  Do not edit: change the order of the
% index and rerun the script (./preview-site.sh and CI do this).
%
% Usage in a sheet, after \newcounter{chap}:
%   % Chapter numbers, generated from 20261125T012000--grh-workbook__guide.org
% by scripts/make-chapters.py.  Do not edit: change the order of the
% index and rerun the script (./preview-site.sh and CI do this).
%
% Usage in a sheet, after \newcounter{chap}:
%   % Chapter numbers, generated from 20261125T012000--grh-workbook__guide.org
% by scripts/make-chapters.py.  Do not edit: change the order of the
% index and rerun the script (./preview-site.sh and CI do this).
%
% Usage in a sheet, after \newcounter{chap}:
%   \input{chapters}\chapnumber{<sheet ID>}
\makeatletter
\providecommand\chapters@set[2]{\@namedef{chapters@#1}{#2}}
\chapters@set{20261125T001000}{1}% zeros-to-ramp-sums
\chapters@set{20261125T001100}{2}% inert-primes-and-estimate-b
\chapters@set{20261125T005000}{3}% galois-fourier-and-averaging
\chapters@set{20261125T005500}{4}% columns-from-a-sphere
\chapters@set{20261125T000000}{5}% slope-method
\chapters@set{20261126T030100}{6}% a-twist-for-each-prime
\chapters@set{20261125T003000}{7}% the-chain-and-the-square-root
\chapters@set{20261125T010000}{8}% rank-lemma-and-zero-estimate
\chapters@set{20261125T010500}{9}% poincare-series-and-near-orthogonality
\chapters@set{20261124T223000}{10}% elliptic-curves-for-the-proof
\chapters@set{20261125T009500}{11}% fock-model-and-theta-intertwiners
\chapters@set{20261124T223100}{12}% complex-multiplication-and-inert-reduction
\chapters@set{20261125T008000}{13}% one-cm-determinant
\chapters@set{20261124T214500}{14}% supersingular-theta-divisibility
\chapters@set{20261124T214600}{15}% the-isotropic-plane
\chapters@set{20261125T013700}{16}% pairings-on-a-formal-group-of-height-two
\chapters@set{20261125T013500}{17}% hecke-nondegeneracy-on-the-tree
\chapters@set{20261126T030200}{18}% the-bound-is-attained
\chapters@set{20261125T009000}{19}% kronecker-limit-and-chowla-selberg
\chapters@set{20261125T002000}{20}% definite-theta-lifts
\chapters@set{20261125T015500}{21}% siegel-weil-in-the-smallest-case
\chapters@set{20261125T002100}{22}% definite-theta-lifts-and-norms
\chapters@set{20261125T016500}{23}% hecke-operators-and-theta-lifts
\chapters@set{20261125T016700}{24}% divisibility-along-the-chain
\chapters@set{20261125T004000}{25}% denominators-and-sha
\chapters@set{20261125T004100}{26}% tate-twists-and-sha
\chapters@set{20261126T020430}{27}% the-minus-part-by-gauss-sums
\chapters@set{20261125T014100}{28}% congruence-modules
\chapters@set{20261125T014500}{29}% eisenstein-congruences-and-selmer-classes
\chapters@set{20261125T014700}{30}% pseudocharacters-and-the-length-bound
\chapters@set{20261125T011000}{31}% galois-representations-black-box
\chapters@set{20261126T030300}{32}% the-data
\chapters@set{20261126T020210}{33}% bessel-poincare-series-and-their-gram-operators
\chapters@set{20261126T030400}{34}% the-seesaw-identity
\chapters@set{20261126T030500}{35}% divisibility-of-the-composite-lift
\chapters@set{20261126T030600}{36}% theorem-a-short-forms
\chapters@set{20261126T030700}{37}% the-determinant-of-a-chain
\chapters@set{20261126T030800}{38}% the-cost-of-one-lift
\chapters@set{20261126T030900}{39}% denominators-of-projections-to-lifts
\chapters@set{20261126T031000}{40}% twists-with-few-congruences
\chapters@set{20261126T031100}{41}% theorem-c-and-the-main-inequality
\providecommand\chapnumber[1]{%
  \@ifundefined{chapters@#1}%
    {\PackageError{chapters}{Sheet #1 is not listed in the index}%
      {List it in 20261125T012000--grh-workbook__guide.org\MessageBreak
       and run scripts/make-chapters.py.}}%
    {\edef\thechap{\@nameuse{chapters@#1}}}}
\makeatother
\chapnumber{<sheet ID>}
\makeatletter
\providecommand\chapters@set[2]{\@namedef{chapters@#1}{#2}}
\chapters@set{20261125T001000}{1}% zeros-to-ramp-sums
\chapters@set{20261125T001100}{2}% inert-primes-and-estimate-b
\chapters@set{20261125T005000}{3}% galois-fourier-and-averaging
\chapters@set{20261125T005500}{4}% columns-from-a-sphere
\chapters@set{20261125T000000}{5}% slope-method
\chapters@set{20261126T030100}{6}% a-twist-for-each-prime
\chapters@set{20261125T003000}{7}% the-chain-and-the-square-root
\chapters@set{20261125T010000}{8}% rank-lemma-and-zero-estimate
\chapters@set{20261125T010500}{9}% poincare-series-and-near-orthogonality
\chapters@set{20261124T223000}{10}% elliptic-curves-for-the-proof
\chapters@set{20261125T009500}{11}% fock-model-and-theta-intertwiners
\chapters@set{20261124T223100}{12}% complex-multiplication-and-inert-reduction
\chapters@set{20261125T008000}{13}% one-cm-determinant
\chapters@set{20261124T214500}{14}% supersingular-theta-divisibility
\chapters@set{20261124T214600}{15}% the-isotropic-plane
\chapters@set{20261125T013700}{16}% pairings-on-a-formal-group-of-height-two
\chapters@set{20261125T013500}{17}% hecke-nondegeneracy-on-the-tree
\chapters@set{20261126T030200}{18}% the-bound-is-attained
\chapters@set{20261125T009000}{19}% kronecker-limit-and-chowla-selberg
\chapters@set{20261125T002000}{20}% definite-theta-lifts
\chapters@set{20261125T015500}{21}% siegel-weil-in-the-smallest-case
\chapters@set{20261125T002100}{22}% definite-theta-lifts-and-norms
\chapters@set{20261125T016500}{23}% hecke-operators-and-theta-lifts
\chapters@set{20261125T016700}{24}% divisibility-along-the-chain
\chapters@set{20261125T004000}{25}% denominators-and-sha
\chapters@set{20261125T004100}{26}% tate-twists-and-sha
\chapters@set{20261126T020430}{27}% the-minus-part-by-gauss-sums
\chapters@set{20261125T014100}{28}% congruence-modules
\chapters@set{20261125T014500}{29}% eisenstein-congruences-and-selmer-classes
\chapters@set{20261125T014700}{30}% pseudocharacters-and-the-length-bound
\chapters@set{20261125T011000}{31}% galois-representations-black-box
\chapters@set{20261126T030300}{32}% the-data
\chapters@set{20261126T020210}{33}% bessel-poincare-series-and-their-gram-operators
\chapters@set{20261126T030400}{34}% the-seesaw-identity
\chapters@set{20261126T030500}{35}% divisibility-of-the-composite-lift
\chapters@set{20261126T030600}{36}% theorem-a-short-forms
\chapters@set{20261126T030700}{37}% the-determinant-of-a-chain
\chapters@set{20261126T030800}{38}% the-cost-of-one-lift
\chapters@set{20261126T030900}{39}% denominators-of-projections-to-lifts
\chapters@set{20261126T031000}{40}% twists-with-few-congruences
\chapters@set{20261126T031100}{41}% theorem-c-and-the-main-inequality
\providecommand\chapnumber[1]{%
  \@ifundefined{chapters@#1}%
    {\PackageError{chapters}{Sheet #1 is not listed in the index}%
      {List it in 20261125T012000--grh-workbook__guide.org\MessageBreak
       and run scripts/make-chapters.py.}}%
    {\edef\thechap{\@nameuse{chapters@#1}}}}
\makeatother
\chapnumber{<sheet ID>}
\makeatletter
\providecommand\chapters@set[2]{\@namedef{chapters@#1}{#2}}
\chapters@set{20261125T001000}{1}% zeros-to-ramp-sums
\chapters@set{20261125T001100}{2}% inert-primes-and-estimate-b
\chapters@set{20261125T005000}{3}% galois-fourier-and-averaging
\chapters@set{20261125T005500}{4}% columns-from-a-sphere
\chapters@set{20261125T000000}{5}% slope-method
\chapters@set{20261126T030100}{6}% a-twist-for-each-prime
\chapters@set{20261125T003000}{7}% the-chain-and-the-square-root
\chapters@set{20261125T010000}{8}% rank-lemma-and-zero-estimate
\chapters@set{20261125T010500}{9}% poincare-series-and-near-orthogonality
\chapters@set{20261124T223000}{10}% elliptic-curves-for-the-proof
\chapters@set{20261125T009500}{11}% fock-model-and-theta-intertwiners
\chapters@set{20261124T223100}{12}% complex-multiplication-and-inert-reduction
\chapters@set{20261125T008000}{13}% one-cm-determinant
\chapters@set{20261124T214500}{14}% supersingular-theta-divisibility
\chapters@set{20261124T214600}{15}% the-isotropic-plane
\chapters@set{20261125T013700}{16}% pairings-on-a-formal-group-of-height-two
\chapters@set{20261125T013500}{17}% hecke-nondegeneracy-on-the-tree
\chapters@set{20261126T030200}{18}% the-bound-is-attained
\chapters@set{20261125T009000}{19}% kronecker-limit-and-chowla-selberg
\chapters@set{20261125T002000}{20}% definite-theta-lifts
\chapters@set{20261125T015500}{21}% siegel-weil-in-the-smallest-case
\chapters@set{20261125T002100}{22}% definite-theta-lifts-and-norms
\chapters@set{20261125T016500}{23}% hecke-operators-and-theta-lifts
\chapters@set{20261125T016700}{24}% divisibility-along-the-chain
\chapters@set{20261125T004000}{25}% denominators-and-sha
\chapters@set{20261125T004100}{26}% tate-twists-and-sha
\chapters@set{20261126T020430}{27}% the-minus-part-by-gauss-sums
\chapters@set{20261125T014100}{28}% congruence-modules
\chapters@set{20261125T014500}{29}% eisenstein-congruences-and-selmer-classes
\chapters@set{20261125T014700}{30}% pseudocharacters-and-the-length-bound
\chapters@set{20261125T011000}{31}% galois-representations-black-box
\chapters@set{20261126T030300}{32}% the-data
\chapters@set{20261126T020210}{33}% bessel-poincare-series-and-their-gram-operators
\chapters@set{20261126T030400}{34}% the-seesaw-identity
\chapters@set{20261126T030500}{35}% divisibility-of-the-composite-lift
\chapters@set{20261126T030600}{36}% theorem-a-short-forms
\chapters@set{20261126T030700}{37}% the-determinant-of-a-chain
\chapters@set{20261126T030800}{38}% the-cost-of-one-lift
\chapters@set{20261126T030900}{39}% denominators-of-projections-to-lifts
\chapters@set{20261126T031000}{40}% twists-with-few-congruences
\chapters@set{20261126T031100}{41}% theorem-c-and-the-main-inequality
\providecommand\chapnumber[1]{%
  \@ifundefined{chapters@#1}%
    {\PackageError{chapters}{Sheet #1 is not listed in the index}%
      {List it in 20261125T012000--grh-workbook__guide.org\MessageBreak
       and run scripts/make-chapters.py.}}%
    {\edef\thechap{\@nameuse{chapters@#1}}}}
\makeatother

\ifcsname chapters@20261126T031100\endcsname
  \chapnumber{20261126T031100}
\else
  \renewcommand{\thechap}{?}
\fi
\renewcommand{\thesection}{\thechap.\arabic{section}}

\theoremstyle{definition}
\newtheorem{problem}{Problem}[section]
\theoremstyle{plain}
\newtheorem{theorem}{Theorem}[section]

\newcommand{\Z}{{\mathbb{Z}}}
\newcommand{\Q}{{\mathbb{Q}}}
\newcommand{\R}{\mathbb R}
\newcommand{\C}{{\mathbb{C}}}
\newcommand{\Oc}{\mathcal O}
\newcommand{\Lc}{\mathcal L}
\newcommand{\main}{\mathrm{main}}
\newcommand{\rem}{\mathrm{rem}}
\newcommand{\hdir}{h_{\mathrm{dir}}}
\newcommand{\hgt}{\operatorname{ht}}
\newcommand{\Gram}{\operatorname{Gram}}
\newcommand{\lean}[1]{\texttt{#1}}
\setlength{\emergencystretch}{3em}

\title{Theorem C and the main inequality}
\author{}
\date{November 26, 2026}

\begin{document}
\maketitle
\refstepcounter{chap}\label{chap}

\paragraph{Prerequisites.}
Hadamard's and Fischer's inequalities for positive definite Hermitian matrices \cite[\S7.8]{HJ}.  Dirichlet \(L\)-functions: \(L(1, \chi) \ne 0\) for \(\chi\) nonprincipal, induced and primitive characters, and the functional equation \cite[Ch.~4, Ch.~5, Ch.~8, Ch.~9]{Dav}.  From Chapter~\ref{c1-chap} the remainder \(R_{d,q}\), written \(R_q\) here for fixed \(d\), the one-sided criterion and the prime number theorem in progressions; from Chapter~\ref{c3-chap} the total remainder; from Chapter~\ref{c5-chap} heights of families; from Chapter~\ref{c6-chap} the theorem on a twist for each prime; from Chapter~\ref{c32-chap} the data; and the statements of Theorems~A and~B (Chapters~\ref{c36-chap} and~\ref{c33-chap}), of the factorization of the chain determinant (Chapter~\ref{c37-chap}), and of the costs of single lifts at a regular family of twists (Chapters~\ref{c38-chap} and~\ref{c40-chap}).  Section~\ref{sec:toy} uses no earlier chapter.

\medskip
Theorems~A and~B supply many independent target forms, with upper bounds for the lengths of their untwisted lifts and for the index of the family.  How small can the covolume and the index of such a family be together?  Theorem~C gives a lower bound for the logarithmic squared covolume plus twice the logarithmic index: the main amount, times the dimension, less a small error.  Comparing the two bounds gives an upper bound for the remainder of the prime sum.  We first do the comparison in a lattice with one twist, then state and prove Theorem~C and the main inequality in the form of the rewritten Lean proof~\cite{GL}, and finally pass to the one-sided bound and to the theorem.

\section{The comparison in a lattice}\label{sec:toy}

Let \(N \ge 1\) be an integer, and give \(\C^d\) the Hermitian form \(\langle v, w \rangle := N\, \bar v^{\mathsf T} w\).  For \(y_1, \dots, y_n \in \Q^d\) write \(\Gram(y) := \det(\langle y_i, y_j \rangle)_{i,j \le n}\), and let \(H := \log N\).

\begin{problem}\label{pr:toy-int}
If \(y_1, \dots, y_d \in \Z^d\) are linearly independent, then \(\Gram(y) \ge N^d\), so \(\log \Gram(y) \ge dH\).  For \(N = 5\), \(y_1 = (1, 2)\), \(y_2 = (2, -1)\) one has \(\Gram(y) = 625\).
\end{problem}

\begin{problem}\label{pr:toy-index}
Let \(y_1, \dots, y_d \in \Q^d\) be linearly independent, and let \(u \in M_d(\Z)\) with \(\det u \ne 0\) be such that \(\sum_k u_{ik} y_k \in \Z^d\) for every \(i\).  If \(\log \Gram(y) \le c\) and \(\log\lvert\det u\rvert \le \delta\), then \(c + 2\delta \ge dH\).  For \(N = 5\), \(y_1 = (\tfrac12, 1)\), \(y_2 = (0, \tfrac13)\), \(u = \mathrm{diag}(2, 3)\), \(c = \log \Gram(y)\) and \(\delta = \log 6\) there is equality.
\end{problem}

Now let the family be shorter than the dimension.  It is completed to a basis by vectors whose lengths are known.

\begin{problem}\label{pr:toy-complete}
Let \(y_1, \dots, y_n\) and \(z_1, \dots, z_{n'}\) in \(\Z^d\) together form a basis of \(\Q^d\), with \(\langle z_i, z_i \rangle \le e^{\beta}\) for every \(i\).  Then \(\log \Gram(y) + n'\beta \ge dH\).  For \(d = 3\), \(N = 2\), \(y_1 = (1, 1, 0)\), \(y_2 = (0, 1, 1)\), \(z_1 = (1, 0, 0)\): \(\Gram(y) = 12\), \(\Gram(y, z) = 8\), and \(8 \le 12 \cdot 2\).
\end{problem}

\begin{problem}\label{pr:toy-steinitz}
Every linearly independent family \(y_1, \dots, y_n \in \Z^d\) is completed to a basis of \(\Q^d\) by \(d - n\) of the standard basis vectors.  Hence \(\log \Gram(y) \ge nH\).
\end{problem}

Completing with \(d - n\) standard basis vectors costs at most \((d - n)H\) in the logarithm of the Gram determinant, where \(H\) is the logarithm of each vector's squared length.

\begin{problem}\label{pr:toy-hadamard}
If \(\langle y_i, y_i \rangle \le e^{u}\) for \(i = 1, \dots, n\), then \(\Gram(y) \le e^{nu}\); if the \(y_i\) are linearly independent, \(\log \Gram(y) \le nu\).
\end{problem}

The comparison itself is an inequality between real numbers.  Suppose that a family of \(n\) vectors in a space of dimension \(d \ge 1\) satisfies a lower bound and two upper bounds,
\[
  d(H - 3F) \le c + 2\delta , \qquad c \le n(H - \rho + F) , \qquad \delta \le dI ,
\]
for real numbers \(H, F, I, \rho, c, \delta\).  The upper exponent per vector is \(H - \rho + F\), including the error \(F\).

\begin{problem}\label{pr:toy-compare}
If moreover \(H + F \ge 0\) and \(d/2 \le n \le d\), then \(n\rho \le d(4F + 2I)\), and \(\rho \le 2(4F + 2I)\) when \(4F + 2I \ge 0\).
\end{problem}

\begin{problem}\label{pr:toy-sign}
Without the hypothesis \(H + F \ge 0\) the conclusion of Problem~\ref{pr:toy-compare} fails: for \(d = 2\), \(n = 1\), \(F = I = 0\), \(H = -1\), \(\rho = 1\), \(c = -2\) and \(\delta = 0\), all three bounds hold, \(d/2 \le n \le d\), and \(n\rho > d(4F + 2I)\).
\end{problem}

In this lattice the main amount sits in the metric.  In the rewritten proof it enters through the height of the family \(\iota \mapsto Z(1)/Z(\theta_\iota)\), where the twist \(\theta_\iota\) is chosen for each pair \((\ell, \iota)\) (Problem~\ref{pr:volume} below).

\section{Theorem C}\label{sec:thmC}

We fix parameters \(q_0, d\) and data for them as in Chapter~\ref{c32-chap}.  Among the conditions on the data, the scales \(x < x_1 < b\) are integers with
\[
  (\log b)^{100} \le x , \qquad x^{10^6} \le x_1 , \qquad x_1^{10^6} \le b , \qquad (\log x_1)^6 \le \log b , \qquad \log\log b \ge 10^{12} ,
\]
the last because \(\log\log b\) is at least \(10^{12}\) times the parameter constant of that chapter, which is at least \(1\).  We also use
\(a = xb\); the number of lifts \(m\) with \(\sqrt x \le m \le 2\sqrt x\); the target rank \(r\), with \(x^5 \le r \le 3x^5\), the source rank \(k = r + 1 - m\), and the step size \(t = 5r\); and the number \(f \ge 1\) of infinite places of \(E\).  Write
\[
  F := fa\log^4 b , \qquad U := fa\log^3 b ,
\]
and \(\main = mk\,\hdir\) for the main amount, where \(\hdir\) is the direction amount; \(\main \ge 0\).  Let \(V_T\) be the space of target forms, \(d_T := \dim_\C V_T\), and \(\mathcal T\) the set of \(m\)-fold twists, with the trivial twist \(1\).  For \(\xi \in \mathcal T\) let \(\Theta_\xi\) be the chain of lifts on \(V_T\), \(Z(\xi) = \det(\Theta_\xi^*\Theta_\xi)\), and, for \(1 \le j \le m\), let \(u_{\xi^j, j}\) be the \(j\)th single lift at the twist \(\xi^j\), restricted to the good forms on the space \(W_j\) of Chapter~\ref{c37-chap}, and \(Z_j(\xi) := \det(u_{\xi^j,j}^* u_{\xi^j,j})\).  Thus \(Z_j(\theta_\iota)\) below is taken at the twist \(\theta_\iota^j\).  Let \(q(\xi)\) be the norm of the modulus of \(\xi\).  The index \(\iota\) runs over pairs of a prime \(\ell\) and an isomorphism \(\iota \colon \C \simeq \bar\Q_\ell\).

\emph{Heights.}  As in Chapter~\ref{c5-chap}, a family \(u = (u_\iota)\) of complex numbers has height at most \(c\), written \(\hgt(u) \le c\), if there are \(A \in \C^\times\), \(\alpha \in \R\) and a function \(h \colon \mathbb N \to \R_{\ge 0}\) such that \(\lvert \sigma A \rvert \le e^\alpha\) for every automorphism \(\sigma\) of \(\C\), \(\lvert \iota(A u_\iota) \rvert_\ell \le \ell^{h(\ell)}\) for every \(\iota\), and \(\alpha + \sum_{\ell \in S} h(\ell)\log\ell \le c\) for every finite set \(S\) of primes.

\begin{problem}\label{pr:ht-add}
If \(\hgt(u) \le c\) and \(\hgt(v) \le c'\), then \(\hgt(uv) \le c + c'\); and \(\hgt(u^n) \le nc\) for every integer \(n \ge 0\).
\end{problem}

\emph{The target lattice.}  The target forms carry a twisted lattice in the sense of Chapter~\ref{c6-chap}, with twists \(\mathcal T\) and lifts \(\Theta_\xi\), all injective: \(\xi\) is allowed at \(\ell\) when \(\gcd(\ell, q(\xi)) = 1\); the lattice \(\Lc_{\xi,\iota}\) consists of the forms in the target of \(\Theta_\xi\) that are integral at \(\iota\); the squared length \(\lVert v \rVert_\sigma^2\) of \(v \in V_T\) at an automorphism \(\sigma\) of \(\C\) is computed from the \(\sigma\)-conjugate coordinates of \(\Theta_1 v\) and a rational Gram matrix; and the volume at \(\xi\) is \(Z(\xi)\).  A family \(y_1, \dots, y_n\) in \(V_T\) has \emph{covolume at most \(c\)} if, for every \(\sigma\), the Gram determinant of the \(\sigma\)-conjugate coordinates of \(\Theta_1 y_1, \dots, \Theta_1 y_n\) is at most \(e^c\) in absolute value.  It has \emph{index at most \(\delta\)} if there is a function \(e \colon \mathbb N \to \R_{\ge 0}\) with \(\sum_{\ell \in S} e(\ell)\log\ell \le \delta\) for every finite set \(S\) of primes and, for every \(\iota\), a matrix \(u\) with entries in \(\Oc_\iota = \{z \in \C : \lvert\iota(z)\rvert_\ell \le 1\}\) and \(\lvert\iota(\det u)\rvert_\ell \ge \ell^{-e(\ell)}\) such that \(\Theta_\xi(\sum_k u_{ik} y_k) \in \Lc_{\xi,\iota}\) for every \(i\) and every \(\xi\) allowed at \(\ell\).

We recall the theorem of Chapter~\ref{c6-chap} in this setting.  Let \(y_1, \dots, y_n, z_1, \dots, z_{n'}\) be a basis of \(V_T\), where \(y\) has covolume at most \(c\) and index at most \(\delta\).  Let \(\theta\) assign to each \(\iota\) a twist \(\theta_\iota\) allowed at its prime \(\ell\), and let \(\gamma \colon \mathcal T \to \C^\times\).  Suppose that, for real numbers \(b_{\mathrm{short}}\), \(b_{\mathrm{denom}}\), \(b_{\mathrm{vol}}\), (i) \(\gamma(\xi)\Theta_\xi z_i \in \Lc_{\xi,\iota}\) for every \(i\), every \(\iota\) and every \(\xi\) allowed at \(\ell\); (ii) \(\lVert z_i \rVert_\sigma^2 \le e^{b_{\mathrm{short}}}\) for every \(i\) and \(\sigma\); (iii) \(\hgt(\lvert\gamma(\theta_\iota)\rvert^{-2}) \le b_{\mathrm{denom}}\); (iv) \(\hgt(Z(1)/Z(\theta_\iota)) \le b_{\mathrm{vol}}\).  Then
\begin{equation}\label{eq:twist}
  0 \le c + 2\delta + n'(b_{\mathrm{short}} + b_{\mathrm{denom}}) + b_{\mathrm{vol}} .
\end{equation}
In the rewritten Lean proof this is \lean{New.TwistedLattice.covolume\_add\_index\_nonneg}, applied to the target lattice; its general statement concerns a twisted lattice on a finite-dimensional complex inner product space with all lifts injective, under the hypotheses above.

Four inputs feed it.  (a) \emph{Auxiliary forms.}  There are an integer \(e\) and \(\gamma \colon \mathcal T \to \C^\times\) with \(\gamma(1) = 1\) and \(\lvert\gamma(\xi)\rvert = q(\xi)^{e/2}\prod_{0 \le i, j < m} T_{ij}(\xi)\), where each \(T_{ij}(\xi)\) is a nonnegative product of \(t\) ratios of absolute values of partial Hecke \(L\)-values, such that \(V_T\) is spanned by forms \(z\) with \(\gamma(\xi)\Theta_\xi z\) integral at \(\iota\) whenever \(\gcd(\ell, q(\xi)) = 1\) and with \(\lVert z \rVert_\sigma^2 \le \exp(2x_1^2b^2)\) for every \(\sigma\).  (b) \emph{Single lifts.}  There is a family \(\theta\) of twists with small ratios such that \(\hgt(Z_j(1)/Z_j(\theta_\iota)) \le d_T(F/m - k\hdir)\) for \(1 \le j \le m\); this combines the cost of one lift (Chapter~\ref{c38-chap}) with a regular family (Chapter~\ref{c40-chap}).  Small ratios means that each \(\theta_\iota\) is allowed at its prime and that, for every \(1 \le j \le m\), every ramified critical pair \(p\) and each sign \(\epsilon \in \{1, -1\}\), the family \(\iota \mapsto \mathrm{ratio}_p(\theta_\iota^j, \epsilon)^{-1}\) of inverse ratios of Hecke \(L\)-values of Chapter~\ref{c40-chap} has height at most \(x_1\).  (c) \emph{Factorization} (Chapter~\ref{c37-chap}).  There is an integer \(e'\) such that for every \(\xi \in \mathcal T\)
\[
  \frac{Z(1)}{Z(\xi)} = \Bigl( q(\xi)^{e'/2} \prod_{1 \le i < j \le m} B_{ij}(\xi) \Bigr)^{-d_T} \prod_{j=1}^m \frac{Z_j(1)}{Z_j(\xi)} ,
\]
where each \(B_{ij}(\xi)\) is a real product of \(t\) ratios of absolute values of partial Hecke \(L\)-values.  (d) \emph{Small ratios.}  If \(\theta\) has small ratios, then for every integer \(e\) both families \(\bigl(q(\theta_\iota)^{e/2}\prod_{0 \le i,j < m} T_{ij}(\theta_\iota)\bigr)^{-1}\) and \(\bigl(q(\theta_\iota)^{e/2}\prod_{1 \le i < j \le m} B_{ij}(\theta_\iota)\bigr)^{-1}\) have height at most \(m^2tx_1\).

\begin{problem}\label{pr:scales}
The data satisfy \(2m^2tx_1 \le b\), \(4(2x_1^2 + b^{-1}) \le b\), \(b \le F\) and \(4/b^2 \le 1/2\).
\end{problem}

\begin{problem}\label{pr:completion}
If \(n \ge (1 - 4/b^2)d_T\) and \(d_T = n + n'\), then \(n'(2x_1^2b^2 + b) \le d_T b\).
\end{problem}

So the allowance for one completing auxiliary form is \(2x_1^2b^2 + b\), and the total allowance for the \(n'\) completing forms is at most \(d_T b\).

\begin{problem}\label{pr:denom}
For \(\gamma\) as in (a) and \(\theta\) as in (b), \(\hgt(\lvert\gamma(\theta_\iota)\rvert^{-2}) \le b\).
\end{problem}

\begin{problem}\label{pr:volume}
For \(\theta\) as in (b), \(\hgt(Z(1)/Z(\theta_\iota)) \le d_T(b + F - \main)\).
\end{problem}

\begin{problem}\label{pr:C}
Let \(y_1, \dots, y_n \in V_T\) be linearly independent, with \(n \ge (1 - 4/b^2)d_T\), covolume at most \(c\) and index at most \(\delta\).  Then \(V_T\) has a basis \(y_1, \dots, y_n, z_1, \dots, z_{n'}\) with each \(z_i\) as in (a), and \(c + 2\delta \ge d_T(\main - 3F)\).
\end{problem}

Problems~\ref{pr:scales}--\ref{pr:C} prove the following.

\begin{theorem}[Theorem C]\label{thm:C}
Let \(y_1, \dots, y_n\) be linearly independent target forms with \((1 - 4/b^2)d_T \le n\), of index at most \(\delta\) and covolume at most \(c\).  Then
\[
  d_T(\main - 3fa\log^4 b) \le c + 2\delta .
\]
\end{theorem}

\paragraph{In the Lean formalization.}  In the rewritten Lean proof~\cite{GL} Theorem~\ref{thm:C} is \lean{New.Data.integral\_forms\_long}, in \lean{grh/src/New/Determinant/Twist/LongForms.lean}.  In the same file the inputs (a), (b) and (c) are \lean{exists\_span\_isAuxiliaryForm\_eq\_top}, \lean{exists\_hasSmallRatios\_heightLe\_liftDeterminant\_div} and \lean{exists\_chainDeterminant\_one\_div\_eq}, Problem~\ref{pr:completion} is \lean{mul\_completion\_cost\_le} and Problem~\ref{pr:volume} is \lean{heightLe\_chainDeterminant\_one\_div}, all in the namespace \lean{New.Data}; the two bounds of (d) are \lean{New.Data.HasSmallRatios.heightLe\_inv\_rpow\_mul\_prod\_turnRatio} there and \lean{New.Data.HasSmallRatios.heightLe\_inv\_rpow\_mul\_prod\_quotientRatio} in \lean{grh/src/New/Determinant/Chain/NormOfChain.lean}; \(T_{ij}\), \(B_{ij}\) and small ratios are \lean{New.Data.turnRatio}, \lean{New.Data.quotientRatio} and \lean{New.Data.HasSmallRatios}.

\section{The main inequality}\label{sec:main}

Let \(\rem\) be the total remainder of Chapter~\ref{c3-chap}.  Theorems~A and~B concern the space \(S\) of ramified source forms, its subset \(S_\Z\) of forms integral at every \(\iota\), a finite group \(\bar G\), and linear maps \(P_g \colon S \to V_T\) for \(g \in \bar G\), the normalized translates of the Poincar\'e map.  We use them in the following form.
\begin{itemize}
\item[(A1)] \(S\) is spanned by the set \(S_{\mathrm{short}}\) of \(\varphi \in S_\Z\) with \(\lVert P_g\varphi \rVert_\sigma^2 \le \exp(\main - m\,\rem + F)\) for every \(g \in \bar G\) and every \(\sigma\).
\item[(A2)] If \(J \subseteq \bar G\) is finite and every \(y_i\) lies in \(P_g(S_\Z)\) for some \(g \in J\), then \(y_1, \dots, y_n\) has index at most \(2\,\# J \dim S \cdot U\).
\item[(B)] \(d_T > 0\), and there is a finite \(J \subseteq \bar G\) with \(\# J \dim S \le 5d_T\log(1 + 2b^2)\) and \(\dim \sum_{g \in J} P_g(S) \ge (1 - 4/b^2)d_T\).
\end{itemize}
The upper bound for the covolume is Hadamard's inequality: if \(\lVert z_i \rVert_\sigma^2 \le e^u\) for \(i = 1, \dots, n\) and every \(\sigma\), then \(z_1, \dots, z_n\) has covolume at most \(nu\).

\begin{problem}\label{pr:family}
For \(J\) as in~(B), the span of \(\bigcup_{g \in J} P_g(S_{\mathrm{short}})\) is \(\sum_{g \in J} P_g(S)\).  Hence that space has a basis \(y_1, \dots, y_n\) with each \(y_i \in P_g(S_{\mathrm{short}})\) for some \(g \in J\), and \((1 - 4/b^2)d_T \le n \le d_T\).
\end{problem}

\begin{problem}\label{pr:upper}
This family has covolume at most \(n(\main - m\,\rem + F)\) and index at most \(d_T \cdot 10\log(1 + 2b^2)\, U\).
\end{problem}

\begin{problem}\label{pr:compare}
\(n\, m\,\rem \le d_T\bigl(4F + 20\log(1 + 2b^2)\,U\bigr)\).
\end{problem}

\begin{problem}\label{pr:errors}
\(n \ge d_T/2\), and \(8F + 40\log(1 + 2b^2)\, U \le fa\log^5 b\).  Hence \(m\,\rem \le fa\log^5 b\).
\end{problem}

\begin{theorem}[Main inequality]\label{thm:main}
\(m\,\rem \le fa\log^5 b\).
\end{theorem}

The errors per dimension are: \(F\) from the size in Theorem~A, \(20\log(1 + 2b^2)\,U\) from the index in Theorems~A and~B (counted twice, as it enters as \(2\delta\)), and \(3F\) from Theorem~C.

\paragraph{In the Lean formalization.}  In the rewritten Lean proof~\cite{GL} Theorem~\ref{thm:main} is \lean{New.star\_inequality}, in \lean{grh/src/New/Reduction/MainInequality.lean}.  The inputs (A1), (A2) and (B) are \lean{span\_short\_integralRamified\_eq\_top}, \lean{indexLe\_normalizedPoincare} and \lean{enough\_forms}, in the namespace \lean{New.Data} of \lean{grh/src/New/Reduction/ShortAndEnough.lean}, and Hadamard's inequality is \lean{New.TwistedLattice.covolumeLe\_of\_normSqAt\_le}.

\section{From the main inequality to the theorem}\label{sec:end}

Let \(q_0 \ge 1\) and \(d\) be the parameters of the data, so that \(18480 \mid q_0\) and \(d\) is a prime with \(d \equiv q_0/2 - 1 \pmod{q_0}\); then \(d\) is odd and prime to \(q_0\).  Recall that, for \(q \ge 1\) prime to \(d\), \(R_q(y) = \sum_{n \le y} \lambda(n)(1/n - 1/y) - \log y/\varphi(q) - c_q\), with \(\lambda(n) = \Lambda(n)(1 - (n/d))\) for \(n \equiv 1 \pmod q\) and \(\lambda(n) = 0\) otherwise, where \((n/d)\) is the Legendre symbol and \(c_q\) is the constant for which \(R_q(y) \to 0\); these are \(R_{d,q}\) and \(c_{d,q}\) of Chapter~\ref{c1-chap}.  When \(d\) varies we write \(R_q^{(d)} = R_{d,q}\).  We use from Chapter~\ref{c3-chap} the identity
\[
  \rem = a\bigl(\#\Phi_b\,(R_{q_0}(x) - R_{q_0}(x_1)) + \#\Phi_a\, R_1(x_1)\bigr)
\]
for the sets \(\Phi_a\) and \(\Phi_b\) of \(a\)-places and \(b\)-places, together with \(f = \#\Phi_a\,\varphi(q_0)\) and \(0 < \#\Phi_a \le \#\Phi_b \log b\) (Chapter~\ref{c32-chap}).

\begin{problem}\label{pr:Rx}
\(R_{q_0}(x) \le \varphi(q_0)\log^6 b/m + \lvert R_{q_0}(x_1)\rvert + \log b\,\lvert R_1(x_1)\rvert\).
\end{problem}

Two further inputs.  By the prime number theorem in progressions, for every prime \(d \ne 2\), every \(q \ge 1\) prime to \(d\) and every integer \(B \ge 0\), one has \(\lvert R_q(y) \rvert \le (\log y)^{-B}\) for all large \(y\) (Chapter~\ref{c1-chap}).  And for \(0 < \varepsilon < 1/100\) and every large integer \(x\) there are data with this \(x\), with \(\log b \le 2x^\varepsilon\), \(m \ge \sqrt x\) and \(\log x_1 \ge x^{\varepsilon/12}\) (Chapter~\ref{c32-chap}).

\begin{problem}\label{pr:PNT}
Let \(\varepsilon > 0\).  For all large integers \(x\) and every real \(y\) with \(\log y \ge x^{\varepsilon/12}\), one has \(\lvert R_{q_0}(y) \rvert \le x^{-2}\) and \(\lvert R_1(y) \rvert \le x^{-2}\).
\end{problem}

\begin{problem}\label{pr:integers}
Let \(0 < \varepsilon < 1/100\).  For all large integers \(x\), \(R_{q_0}(x) \le (64\varphi(q_0) + 3)\, x^{-1/2 + 6\varepsilon}\).
\end{problem}

\begin{problem}\label{pr:floor}
For \(q \ge 1\) prime to \(d\) and real \(y \ge 2\), \(R_q(y) - R_q(\lfloor y \rfloor) \le 3/y\).  For \(-1 \le s \le 0\), \(\lfloor y \rfloor^{s} \le 2y^{s}\).
\end{problem}

The exercise sharpens the bound \(26/y\) used in the rewritten proof~\cite{GL}, whose lemma \lean{New.primeRemainder\_sub\_primeRemainder\_floor\_le} gives that bound for any prime \(d\), positive modulus \(q\) and real \(y \ge 2\).

\begin{problem}\label{pr:reals}
Let \(\varepsilon > 0\).  For all large real \(y\), \(R_{q_0}(y) \le y^{-1/2 + \varepsilon}\).
\end{problem}

\begin{theorem}[One-sided bound]\label{thm:one-sided}
Let \(q_0 \ge 1\) be a multiple of \(18480\), let \(d \equiv q_0/2 - 1 \pmod{q_0}\) be a prime, and let \(\varepsilon > 0\).  Then \(R_{q_0}(y) \le y^{-1/2 + \varepsilon}\) for all large real \(y\).
\end{theorem}

The criterion of Chapter~\ref{c1-chap} reads: let \(q_0 \ge 1\) be a multiple of \(8\), let \(d_0 \ge 0\) be an integer with \(d_0 \equiv 7 \pmod 8\) and prime to \(q_0\), and let \(\beta > \frac12\); if for every prime \(d \equiv d_0 \pmod{q_0}\) one has \(R^{(d)}_{q_0}(y) \le y^{\beta - 1}\) for all large \(y\), then \(L(\rho, \psi) \ne 0\) for every character \(\psi\) modulo \(q_0\) and every \(\rho \ne 1\) with \(\operatorname{Re}\rho > \beta\).

\begin{problem}\label{pr:d0}
For \(q \ge 1\), \(q_0 := 18480q\) and \(d_0 := 9240q - 1\): \(d_0 \equiv 7 \pmod 8\), \(\gcd(d_0, q_0) = 1\), and \(d_0 = q_0/2 - 1\).
\end{problem}

\begin{problem}\label{pr:induced}
Let \(q \ge 1\), let \(\chi\) be a character modulo \(q\), and let \(\chi'\) be the character modulo \(q_0 = 18480q\) that it induces.  If \(\chi\) is nonprincipal or \(s \ne 1\), then \(L(s, \chi') = L(s, \chi)\prod_{p \mid q_0}(1 - \chi(p)p^{-s})\), so \(L(s, \chi') \ne 0\) implies \(L(s, \chi) \ne 0\).
\end{problem}

Problems~\ref{pr:d0} and~\ref{pr:induced}, Theorem~\ref{thm:one-sided}, the criterion and \(L(1, \chi) \ne 0\) for \(\chi\) nonprincipal prove the following.

\begin{theorem}\label{thm:grh}
Let \(q \ge 1\), let \(\chi\) be a Dirichlet character modulo \(q\), and let \(s \in \C\) with \(\operatorname{Re} s > \frac12\), where \(s \ne 1\) if \(\chi\) is principal.  Then \(L(s, \chi) \ne 0\).
\end{theorem}

\begin{problem}\label{pr:line}
Every zero \(s\) of a Dirichlet \(L\)-function with \(0 < \operatorname{Re} s < 1\) has \(\operatorname{Re} s = \frac12\).
\end{problem}

\paragraph{In the Lean formalization.}  In the rewritten Lean proof~\cite{GL}, Problem~\ref{pr:Rx} is \lean{New.Data.primeRemainder\_le}, Theorem~\ref{thm:one-sided} is \lean{New.eventually\_primeRemainder\_le\_rpow} and Theorem~\ref{thm:grh} is \lean{New.grh}, all in \lean{grh/src/New/Reduction/OneSidedBound.lean}; the deduction of Problem~\ref{pr:line} from Theorem~\ref{thm:grh} is \lean{New.zeros\_on\_line}, in \lean{grh/src/New/Main.lean}.  The identity for \(\rem\) is \lean{New.Data.totalRemainder\_eq}, the prime number theorem in progressions is \lean{New.eventually\_abs\_primeRemainder\_le}, the existence of data is \lean{New.eventually\_exists\_data\_scales}, the criterion is \lean{New.LFunction\_ne\_zero\_of\_primeRemainder\_le}, and Problems~\ref{pr:PNT} and~\ref{pr:integers} are \lean{New.eventually\_abs\_primeRemainder\_exp\_le} and \lean{New.eventually\_primeRemainder\_natCast\_le}.

\section*{Hints}

\subsection*{\S\ref{sec:toy}}
\begin{description}
\item[\ref{pr:toy-int}] \((\langle y_i, y_j \rangle) = N\, Y Y^{\mathsf T}\) for the matrix \(Y\) with rows \(y_i\), and \(\det Y\) is a nonzero integer.
\item[\ref{pr:toy-index}] Problem~\ref{pr:toy-int} for the rows of \(uY\).
\item[\ref{pr:toy-complete}] Fischer's inequality for the Gram matrix of \((y, z)\), Hadamard's for \(\Gram(z)\), then Problem~\ref{pr:toy-int}.
\item[\ref{pr:toy-steinitz}] Steinitz exchange.
\item[\ref{pr:toy-hadamard}] Apply Hadamard's inequality to the Gram matrix.
\item[\ref{pr:toy-compare}] Insert the upper bounds into the lower bound; then use \(n(H + F) \le d(H + F)\).
\end{description}

\subsection*{\S\ref{sec:thmC}}
\begin{description}
\item[\ref{pr:ht-add}] Multiply the numbers \(A\), and add the \(\alpha\) and the \(h(\ell)\).
\item[\ref{pr:scales}] \(m^2 \le 4x\), \(t \le 15x^5\) and \(x \ge 2\).
\item[\ref{pr:completion}] \(n' \le 4d_T/b^2\).
\item[\ref{pr:denom}] Input~(d) and Problems~\ref{pr:ht-add} and~\ref{pr:scales}.
\item[\ref{pr:volume}] Input~(c); the scalar by~(d), the single lifts by~(b).
\item[\ref{pr:C}] Apply~\eqref{eq:twist} with \(b_{\mathrm{short}} = 2x_1^2b^2\), \(b_{\mathrm{denom}} = b\) and \(b_{\mathrm{vol}} = d_T(b + F - \main)\).
\end{description}

\subsection*{\S\ref{sec:main}}
\begin{description}
\item[\ref{pr:family}] Images of a spanning set span the image.
\item[\ref{pr:upper}] Hadamard's inequality with \(u = \main - m\,\rem + F\); (A2) and the bound for \(\# J\dim S\) in~(B).
\item[\ref{pr:compare}] Compare the two bounds for \(c + 2\delta\), using \(n \le d_T\) and \(\main + F \ge 0\).
\item[\ref{pr:errors}] \(\log(1 + 2b^2) \le 3\log b\), and \(\log b \ge 128\).
\end{description}

\subsection*{\S\ref{sec:end}}
\begin{description}
\item[\ref{pr:Rx}] Solve the identity for \(R_{q_0}(x)\), and divide Theorem~\ref{thm:main} by \(ma\#\Phi_b\).
\item[\ref{pr:PNT}] Take \(B \ge 24/\varepsilon\); \(d\) is odd and prime to \(q_0\).
\item[\ref{pr:integers}] Problems~\ref{pr:Rx} and~\ref{pr:PNT} at \(y = x_1\).
\item[\ref{pr:floor}] The two sums differ by \((1/\lfloor y \rfloor - 1/y)\sum_{n \le y}\lambda(n)\); use \(0 \le \lambda \le 2\Lambda\) and \(\psi(\lfloor y \rfloor) < 1.04\lfloor y \rfloor\) \cite[Thm.~12]{RS}.
\item[\ref{pr:reals}] Problem~\ref{pr:integers} with \(\min(\varepsilon/7, 1/200)\) in place of \(\varepsilon\), and Problem~\ref{pr:floor}.
\item[\ref{pr:induced}] Compare Euler products for \(\operatorname{Re} s > 1\), then continue analytically.
\item[\ref{pr:line}] Reduce to a primitive character; then the functional equation and Theorem~\ref{thm:grh}.
\end{description}

\begin{thebibliography}{GL}
\bibitem[Dav]{Dav} H.~Davenport, \emph{Multiplicative Number Theory}, 3rd ed., Grad.\ Texts in Math.\ 74, Springer, 2000.
\bibitem[GL]{GL} L.~Sernau, \emph{GRH for Dirichlet \(L\)-functions, in Lean~4}, a shorter rewriting of the Lean proof, text by Claude, \url{https://github.com/lsernau/grh-lean}, commit \texttt{72fead1}, readable edition \texttt{grh/}.
\bibitem[HJ]{HJ} R.~A.~Horn and C.~R.~Johnson, \emph{Matrix Analysis}, 2nd ed., Cambridge Univ.\ Press, 2013.
\bibitem[RS]{RS} J.~B.~Rosser and L.~Schoenfeld, \emph{Approximate formulas for some functions of prime numbers}, Illinois J.\ Math.\ \textbf{6} (1962), 64--94.
\end{thebibliography}

\end{document}
